\scanpage{75}
\mspage{37}
d'inertie $I_x=T_\infty(\overline K)$, cf plus haut. Soit $f_K$ $\unsure{K}$ un pt de $U_K$ rat/$K$, d'où une section $\psi=\psi_{f_K}$ de $E_K$ sur $E_{U_K}$. Ceci posé les conditions suivantes sont équivalentes

\noindent a) $f_K$ se prolonge en une section de $U$ sur $S$

\noindent b) $\psi$ provient d'une section de $E_U$ sur $E_S$

\noindent b') le composé $E_K\xrightarrow{\psi}E_{U_K}\longrightarrow E_U$ s'annule sur $I_x$

\deleted{\unsure{c)}}

L'équivalence de b) et b'), et qu'ils soient impliqués par a), est clair. C'est l'implication $\mathrm{b)}\Longrightarrow\mathrm{a)}$ qui demande une démonstration. On est ramené aussitôt, par descente, au cas \deleted{\unsure{d'un}} où $T$ est strictement local (donc $E_K=\operatorname{Gal}\overline K/K$ est réduit à son ss-groupe d'inertie, \inserted{et $E_S=(1)$}). On est ramené de proche en proche au cas où $U/S$ est une courbe relative élémentaire, \deleted{de genre $g\geq 1$,} \inserted{$U=X\setminus T$, $X$ lisse et propre.} Alors $f$ se prolonge en une section $f$ \deleted{rationnelle}\inserted{de} $X$ sur $S$, et la conclusion est que $f(S)\subset U$. Donc on est ramené finalement au

\uline{Lemme} Soit $X$ schéma projectif lisse de dim relative 1 connexe sur $\unsure{S|T}$ trait strictement local,
\margin{74}
